\section{Ballot, convolution, rational, Fuss, and super-Catalan extensions}
\label{sec:generalizations}

\subsection{Bertrand's ballot theorem at excess one}

Consider words with $n+1$ votes $A$ and $n$ votes $B$.  Give $A$ weight $+1$
and $B$ weight $-1$.  The cycle lemma says that a word with total weight
$r>0$ has exactly $r$ distinguished cyclic starting positions whose nonempty
partial sums are positive, counted with positions even when the underlying
word is periodic.  For total weight one, every pointed word has exactly one
good start.  Double counting pointed words gives
\[
 \frac1{2n+1}\binom{2n+1}{n}=\Cat_n.                             \tag{7.1}
\]
Removing the first $A$ from a positive word gives a Dyck word.  This proves the
ballot interpretation and Equation~\eqref{eq:odd-binomial} simultaneously.

For \ipra, the cycle lemma can be proved by the deletion algorithm mentioned
on Wikipedia: place the letters on a circle, repeatedly delete adjacent
$AB$ pairs, and record the surviving $A$ positions.  The deletion relation is
bounded, a fixed least-pair strategy makes it primitive recursive, and an
induction on word length proves that the survivors are exactly the good
starts.  For total weight $r$, exactly $r$ survive.

\subsection{The \(k\)-fold Catalan convolution}

A forest of $k$ ordered full binary trees with a total of $n$ internal nodes is
counted in two ways.  Splitting by the internal-node counts of the components
gives
\[
 \sum_{i_1+\cdots+i_k=n}\Cat_{i_1}\cdots\Cat_{i_k}.              \tag{7.2}
\]
Encode the forest in Polish order, writing an internal node with weight $+1$
and a leaf with weight $-1$.  There are $n$ internal nodes and $n+k$ leaves,
so the word length is $2n+k$ and its total weight is $-k$.  The Raney form of
the cycle lemma says that, among the $2n+k$ pointed rotations of every such
word, exactly $k$ have all proper prefix sums greater than $-k$ in the order
needed to cut the word uniquely into $k$ trees.  Hence
\[
 \boxed{\displaystyle
 \sum_{i_1+\cdots+i_k=n}\Cat_{i_1}\cdots\Cat_{i_k}
  =\frac{k}{2n+k}\binom{2n+k}{n}.}                               \tag{7.3}\label{eq:kfold}
\]
The internal theorem is
$(2n+k)$ times the left side equals $k\binom{2n+k}{n}$.
The pointed-word proof avoids any assumption that cyclic orbits have full
size; periodic words cause no difficulty because the distinguished position is
part of the code.

Coefficientwise, Equation~\eqref{eq:kfold} is also Lagrange inversion for
$C(x)^k$.  The forest proof is the smaller \ipra{} dependency.

\subsection{Rational Catalan numbers}

Let $r,s$ be coprime positive integers.  Give a right step weight $s$ and an up
step weight $-r$.  A path from $(0,0)$ to $(r,s)$ has total weight zero.  Mark
one step and perform the standard cycle-lemma normalization at the least
partial sum.  Coprimality implies that the $r+s$ cyclic partial-sum ranks are
distinct modulo $r+s$, so exactly one marked rotation stays on the required
side of the line $ry=sx$.  Double counting gives
\[
 \operatorname{RatCat}(r,s)
 =\frac1{r+s}\binom{r+s}{r}.                                    \tag{7.4}\label{eq:rational-cat}
\]
The proof is constructive: compute all $r+s$ partial sums, choose the least
rank with the fixed tie convention, rotate, and invert by remembering the
mark.  Coprimality is used through Bezout data obtained by the primitive-
recursive extended Euclidean algorithm.  Equation~\eqref{eq:rational-cat}
is represented after multiplication by $r+s$.

The ordinary Catalan case is $(r,s)=(n,n+1)$ after the strict/weak boundary
translation; Equation~\eqref{eq:odd-binomial} is the resulting formula.

\subsection{Fuss--Catalan numbers}

For $p\ge2$ and $k\ge0$, define
\[
 \Cat^{(p)}_0=1,
 \qquad
 \Cat^{(p)}_k
 =\frac1{(p-1)k+1}\binom{pk}{k}
 =\frac1k\binom{pk}{k-1}\quad(k\ge1).                           \tag{7.5}\label{eq:fuss}
\]
The two fractions are equal by adjacent-binomial arithmetic.

A full ordered $p$-ary tree with $k$ internal nodes has $(p-1)k+1$ leaves.
Its Polish word therefore has $k$ symbols of weight $p-1$ and
$(p-1)k+1$ symbols of weight $-1$, total length $pk+1$ and total weight $-1$.
The cycle lemma gives exactly one valid pointed rotation.  Thus the number of
trees is
\[
 \frac1{pk+1}\binom{pk+1}{k}
 =\frac1{(p-1)k+1}\binom{pk}{k},
\]
proving Equation~\eqref{eq:fuss}.  Recursive parsing of the Polish word is the
primitive-recursive inverse.

The same trees are syntax trees for $k$ applications of a fixed $p$-ary
operator.  Rooting a dissection at a boundary edge and taking the incident
$(p+1)$-gon gives a $p$-tuple of smaller dissections, so they also count
partitions of a convex $((p-1)k+2)$-gon into $k$ nonoverlapping
$(p+1)$-gons.  The Lukasiewicz height of the Polish word gives the $p$-good
path interpretation.  These are explicit recursive bijections, not merely
recurrence comparisons.

The formal series $T=1+xT^p$ has coefficients $\Cat^{(p)}_k$.  For every $q\ge1$, finite Lagrange inversion gives
\[
 [x^k]T^q=\frac{q}{pk+q}\binom{pk+q}{k},                         \tag{7.6}\label{eq:fuss-power}
\]
with the denominator cleared.  Forests of $q$ $p$-ary trees give the same
identity by the Raney lemma.

\subsection{The two-parameter path numbers on MathWorld}

The notation on MathWorld can be handled uniformly by formal reversion.  Fix
$p>1$ and $0\le q\le p-1$.  Let $X=X(y)$ be the unique formal series with
\[
 X=y(1+X)^p.                                                      \tag{7.7}\label{eq:Xreversion}
\]
Define
\[
 D_{p,q}(y)=\frac{X}{(1+X)^q}=\sum_{k\ge1}d_{p,q}(k)y^k.          \tag{7.8}\label{eq:Dpq}
\]
Finite Lagrange inversion gives
\[
 d_{p,q}(k)=\frac{p-q}{pk-q}\binom{pk-q}{k-1}.                  \tag{7.9}\label{eq:Dexplicit}
\]
For $q=0$ this is the Fuss--Catalan number in
Equation~\eqref{eq:fuss}.  The coefficient proof is a bounded sum: from
\([y^k]\Phi(X)=k^{-1}[t^{k-1}]\Phi'(t)(1+t)^{pk}\), take
$\Phi(t)=t(1+t)^{-q}$ and simplify by Pascal's identity.  Negative exponents
are avoided by multiplying by a sufficiently large power of $1+t$ before
coefficient extraction.

Substituting $y=x/(1+x)^p$ in Equation~\eqref{eq:Dpq} gives
\[
 D_{p,q}\left(\frac{x}{(1+x)^p}\right)=\frac{x}{(1+x)^q}.        \tag{7.10}\label{eq:Dsub}
\]
Taking the coefficient of $x^k$ after multiplication by $(1+x)^n$ proves the
generalized Jonah identity displayed by MathWorld:
\[
 \binom{n-q}{k-1}
 =\sum_{i=1}^{k}d_{p,q}(i)\binom{n-pi}{k-i}.                    \tag{7.11}\label{eq:jonah}
\]
The natural generalized-binomial convention applies if an upper entry is
negative; in the ordinary path range all terms are nonnegative and finite.

For $r>0$ with $q<p-r$, Equation~\eqref{eq:Dpq} also gives
\[
 D_{p,q}(y)=\frac{D_{p,p-r}(y)D_{p,q+r}(y)}{y}.                  \tag{7.12}
\]
Comparing the coefficient of $y^k$ yields the MathWorld recurrence
\[
 d_{p,q}(k)=
 \sum_{\substack{i,j>0\\i+j=k+1}}
 d_{p,p-r}(i)d_{p,q+r}(j).                                      \tag{7.13}\label{eq:Dconvolution}
\]
Equations~\eqref{eq:Dexplicit}, \eqref{eq:jonah}, and
\eqref{eq:Dconvolution} therefore have a single coefficientwise \ipra{} proof.
The path interpretation follows by decomposing at the first contact with the
appropriate $p$-slope boundary; the formal proof is independent of drawing
conventions.

\subsection{Super-Catalan integers}

The two-parameter super-Catalan numbers on Wikipedia are
\[
 S(m,n)=\frac{(2m)!(2n)!}{m!n!(m+n)!}.                           \tag{7.14}\label{eq:supercat}
\]
Their integrality is not obvious from the quotient, but it has a short
valuation proof suitable for \ipra.  For every prime power $a=p^j$,
\[
 \floor{\frac{2m}{a}}+\floor{\frac{2n}{a}}
 \ge
 \floor{\frac m a}+\floor{\frac n a}+\floor{\frac{m+n}{a}}.     \tag{7.15}\label{eq:super-floor}
\]
After removing integer parts, this is
$\floor{2u}+\floor{2v}\ge\floor{u+v}$ for $0\le u,v<1$, verified by
the four cases $u<1/2$ or $u\ge1/2$, and similarly for $v$.  Summing
Equation~\eqref{eq:super-floor} over $j$ with Legendre's formula shows
$\vp p(S(m,n))\ge0$ for every prime $p$, hence the denominator divides the
numerator.  In bounded arithmetic, a least-prime-divisor induction supplies
the final divisibility conclusion.

Direct cancellation gives
\[
 S(1,n)=2\Cat_n,
 \qquad
 S(n,n)=\binom{2n}{n}.                                           \tag{7.16}
\]
The source correctly notes that a uniform combinatorial interpretation for
all $(m,n)$ is open.  This report proves integrality and the displayed special
cases but does not manufacture a nonexistent general bijection.

\subsection{Coxeter Catalan statements: type \(A\) and two source corrections}

For the standard generalized Catalan number of the rank-$n$ root system
$A_n$ \cite{FominReading}, identify a positive root with an interval $[i,j]$.  Order ideals in the
triangular interval poset are in primitive-recursive bijection with Dyck paths:
trace the northeast boundary of the ideal.  Taking maximal elements and taking
downward closure are inverse bounded maps between order ideals and antichains.
Consequently
\[
              \operatorname{Cat}(A_n)=\Cat_{n+1}.               \tag{7.17}
\]
Thus, in the standard rank convention, the ordinary number $\Cat_n$
corresponds to $A_{n-1}$, not $A_n$.

Two sentences on the requested Wikipedia revision need correction.
First, root-poset antichains are not equinumerous with fully commutative group
elements in every finite crystallographic type \cite{StembridgeFC}.  The smallest counterexample
is $B_2=I_2(4)$.  Its four positive roots have the two simple roots incomparable below $\alpha_1+\alpha_2$, which is below the highest root; their antichains are
\[
 \varnothing,\quad \{\alpha_1\},\quad \{\alpha_2\},\quad
 \{\alpha_1+\alpha_2\},\quad \{\alpha_1+2\alpha_2\},\quad
 \{\alpha_1,\alpha_2\},
\]
so there are $6$.  The dihedral Coxeter group has $8$ elements, and every
one except the longest element is fully commutative, so there are $7$.
The two counts do agree in type $A$: fully commutative permutations are the
$321$-avoiding permutations already counted in Section~\ref{sec:combinatorics}.

Second, the unweighted assertion that a standard generalized Catalan number
is the sum over all maximal proper subdiagrams is false.  For $A_2$ the left
side is $\Cat_3=5$, whereas the two deleted diagrams are $A_1$, each with
Catalan number $\Cat_2=2$, giving $4$.  In type $A$ the correct weighted
maximal-parabolic recurrence is
\[
 2n\,\operatorname{Cat}(A_n)
  =(n+3)\sum_{i=1}^{n}
      \operatorname{Cat}(A_{i-1})\operatorname{Cat}(A_{n-i}).  \tag{7.18}\label{eq:coxeter-A-recurrence}
\]
Indeed, with $\operatorname{Cat}(A_j)=\Cat_{j+1}$ and
$\operatorname{Cat}(A_0)=1$, the sum is
\[
 S_n=\sum_{i=1}^{n}\Cat_i\Cat_{n+1-i}
     =\Cat_{n+2}-2\Cat_{n+1}
     =\frac{2n}{n+3}\Cat_{n+1},
\]
where the middle equality is Segner and the last is the ratio recurrence.
Clearing $n+3$ gives Equation~\eqref{eq:coxeter-A-recurrence} as a finite
\ipra{} identity.  The unweighted Segner sum does hold for the shifted
``positive'' type-$A$ Catalan invariant $\operatorname{Cat}^{+}(A_n)=\Cat_n$;
this likely explains the conflation on the source page.

The general root-poset theorem and the correctly normalized parabolic
recurrence for every finite crystallographic type require the classification
of finite root systems (or a uniform Coxeter-theoretic development).  They are
recorded as formalization targets, not asserted as consequences of the false
unweighted sentence.
